% 导言区：在 commons/cumcmthesis 的基础上扩展，不直接改 cls。

% 图片目录（路径 图片/图X.png）
\graphicspath{{图片/}}

% 中文主字体设为宋体（SimSun），并启用伪粗体。
% 关键在于把粗体形状挂在“宋体”上，\bfseries/\textbf 才会渲染成加粗宋体，
% 而不是回落到默认的粗体替换字体（黑体）。
\setCJKmainfont[AutoFakeBold={4}]{SimSun}
% 同时把“宋体”族也绑定同一伪粗体，确保 \songti \song 与加粗同时使用时也是加粗宋体
\setCJKfamilyfont{song}{SimSun}[AutoFakeBold={4}]
\setCJKfamilyfont{zhsong}{SimSun}[AutoFakeBold={4}]
% 公式中的汉字（如“宋体”里的下标汉字）需要以下两行：
%   CJKmath=true 允许在数学模式里直接写汉字；\setCJKmathfont 指定公式里的汉字字体
\XeTeXlinebreaklocale "zh"
\xeCJKsetup{CJKmath=true}
\setCJKmathfont{SimSun}

% 附录代码排版：listings 直接引入 支持文件 目录下的 .py 源码。
% xeCJK 会处理列表中的中文注释。
\usepackage{listings}
\lstset{
  language=Python,
  basicstyle=\zihao{6},
  breaklines=true,
  columns=flexible,
  keepspaces=true,
  showstringspaces=false,
  keywordstyle=\bfseries,
  tabsize=4,
  aboveskip=6pt,
  belowskip=10pt,
}
